% =====================================================================
% Section III -- Background: the UVM fault path and its prefetcher
% SpecAsync-UVM manuscript, draft 1 (2026-10-10)
%
% SCOPE: the STOCK driver only. SpecAsync's design (worker, queue, width,
% oracle) belongs to Section IV; results belong to Sections VI-VII. This
% section states only what the source and earlier gates establish, so that
% Section I's mechanism sentences have a place to point to.
%
% SOURCES (function names are cited in prose; line numbers stay here):
%   driver 595.91.07, kernel-open/nvidia-uvm (pristine DKMS copy):
%     uvm_perf_prefetch.c  :42 default 51; :48 declaration; :60 module_param
%       (0444); :94-100 subregion count; :102-146 compute_prefetch_region;
%       :118 strict test counter*100 > subregion_pages*threshold; :148-219
%       grow_fault_granularity (64 KB big-page fill, density bitmap only);
%       :226-227 bitmap = resident_mask | faulted_pages; :396-397 destination
%       resident mask; :404-407 first-touch whole-block rule; :552-561 range
%       check (>100 silently -> 51)                        [E5_MECHANISM]
%     uvm_va_block.c :11487 get_prefetch_hint; :11520 PREFETCH access type;
%       :11532 prefetch mask OR'd into new_residency; :11926 service_finish;
%       :11947 did_not_migrate (not removed); :11961 protection loop over all
%       of new_residency; :11770/:11981 block_service_finish_map, the only
%       mapping call on the fault path; :4968-4995 make_resident = copy +
%       finish, no mapping; :4886-4887 "any move implies mappings removed"
%                                              [E5_MECHANISM, E0.9a-2 §3]
%     uvm_migrate.c :680 do_mappings for GPU destination; :698-699
%       MAKE_RESIDENT_AND_MAP; :265-268/:152 block_migrate_add_mappings;
%       :94-101 uvm_va_block_map RWA            [PREFETCH_ASYNC_MAPPING.md]
%   D5 classification (CPU orchestration + async push via uvm_push_end, the
%     only tracker wait gated on CPU destination): D5_CHARACTERIZATION.md,
%     traced on 595.71.05; the classifying functions are byte-identical to
%     v580.95.05; NOT re-traced on 595.91.07 -> [SRC-1].
%   Phase shares: CS2-1 (63-74% T4; 52.1-76.9% 5070 Ti 595.84; UNCHECKED in
%     C1, carried), CS2-3 (D3 <0.2%; 0.05-0.14% on 5070 Ti).
%
% OPEN:
%   [SRC-1] re-confirm the D5 classification on 595.91.07 (read-only diff of
%           service_copy / make_resident_copy / block_copy_end_push).
%   [SRC-2] eviction granularity and policy (LRU over 2 MB root chunks) is
%           cited from [go2023earlyadaptor, ganguly2019interplay], not read
%           from the 595.91.07 source. Read uvm_pmm_gpu.c before submission.
%   [SRC-3] whether CPU-destination faults run the same density rule with the
%           same threshold (E9's CPU-fault counts suggest so; C4 step 5).
%           The text below makes NO claim about it.
%   [FIG]   Fig. 1 (fault path + where each mechanism acts) to be drawn.
% =====================================================================

\section{Background: The UVM Fault Path and Its Prefetcher}
\label{sec:background}

This section describes the parts of NVIDIA's open-source UVM driver that
decide our results: how a GPU page fault is serviced, the distinction between
making a page \emph{resident} and making it \emph{mapped}, and the in-path
prefetcher's density rule. Everything here refers to the stock driver,
release 595.91.07 unless stated otherwise; function names are those of that
release.

\subsection{Servicing a Replayable Fault}
\label{sec:bg-fault-path}

UVM manages memory in \emph{VA blocks}: aligned 2\,MB regions of virtual
address space, each with its own lock and its own per-processor state. When a
GPU warp touches a page that is not mapped in the GPU's page tables, the GPU
records a replayable fault in a hardware fault buffer and the faulting warps
stall. A driver bottom half drains that buffer and services the faults in
\emph{batches}, which prior measurement identified as the basic unit of UVM
fault-handling work \cite{allen2021indepth,allen2024finegrain}.

We decompose one batch into the phases of Table~\ref{tab:phases}, measured by
instrumentation added at their boundaries (Section~V). The batch is fetched
(D1) and preprocessed (D2): entries are sorted, duplicates removed, and faults
to the same page coalesced. The driver then takes the address-space lock in
read mode (D3 is the wait, D4 the hold) and, inside the hold, dispatches the
batch block by block (D5). For each VA block it takes the block's lock,
classifies the faults, chooses where each page should reside, asks the
prefetcher for a hint, migrates whatever is not yet at its destination, and
finally maps the pages for the faulting GPU. Once all blocks are serviced the
driver tells the GPU to replay the stalled accesses (D6).

\begin{table}[t]
\centering
\footnotesize
\setlength{\tabcolsep}{4pt}
\begin{tabular}{@{}l p{0.52\columnwidth} l@{}}
\toprule
\textbf{Phase} & \textbf{Work} & \textbf{Lock} \\
\midrule
D1 & Drain the fault buffer & --- \\
D2 & Sort, de-duplicate, coalesce & --- \\
D3 & Wait for the address-space lock & --- \\
D4 & Hold the address-space lock (contains D5) & read \\
D5 & Per-block dispatch: classify, select residency, prefetch hint,
     copy, map & read + block \\
D6 & Push the replay to the GPU & --- \\
D7 & Residual & --- \\
\bottomrule
\end{tabular}
\caption{Phases of one fault batch, as instrumented in this study.}
\label{tab:phases}
\end{table}

Two properties of this path matter later. First, the lock is rarely
contended in our single-stream workloads: lock \emph{wait} (D3) is below
0.2\% of the batch in the typical case on both of our platforms, while the
lock \emph{hold}, dominated by D5, is the largest phase (63--74\% of the
batch window on the T4, 52--77\% on the RTX~5070~Ti with driver 595.84).
Second, D5 is CPU work, not a wait for the GPU. Its cost is fault
classification, residency bookkeeping, and the construction of copy commands,
which are submitted asynchronously and tracked rather than waited on; the one
blocking wait in this call chain applies only when the destination is the CPU,
which is never the case for the GPU faults we measure. Moving a fault's
discovery earlier does not, by itself, remove this work.

\subsection{Residency Is Not Mapping}
\label{sec:bg-resident-mapped}

UVM tracks two separate facts about every page and processor. A page is
\emph{resident} on a processor when that processor's memory holds the current
copy of its data. It is \emph{mapped} for that processor when the processor's
page tables contain a valid entry for it. A GPU access faults on the second
condition, not the first: a page can be resident in GPU memory and still fault
if the GPU has no entry for it.

The driver changes these two states through different functions. Its
\texttt{make\_resident} routine copies data to the destination and updates the
residency masks; it installs no page-table entry, and its own documentation
states that moving a page implies its existing GPU mappings have been removed.
Mappings for faulting accesses are installed at one place on the fault path,
the mapping step at the end of block service
(\texttt{block\_service\_finish\_map}). Explicit migrations requested from
user space take a third route: a prefetch to a GPU through the driver's
migration interface runs in a make-resident-\emph{and}-map mode and installs
read-write-atomic mappings for the destination GPU, so a page prefetched that
way does not fault when first touched.

This distinction is the hinge of our study. A mechanism that only makes pages
resident ahead of demand can make the later fault \emph{cheaper}, since the
copy inside D5 becomes a no-op, but it cannot make the fault disappear. Only
a mapping step can do that, and on the fault path the mapping step belongs to
the servicing thread.

\subsection{The Density Prefetcher}
\label{sec:bg-prefetcher}

The servicing thread does not map only the faulting pages. Before migrating,
it consults an in-path prefetcher that may enlarge the set of pages to be
serviced. The prefetcher descends from the tree-based neighbourhood scheme
studied by Ganguly \emph{et al.}\ \cite{ganguly2019interplay} and works within
one VA block at a time.

For a fault on a GPU, the prefetcher builds a bitmap over the block's pages.
A page is set if it is \emph{already resident on the destination GPU} or is
being faulted in this batch; faulted pages are first widened to their 64\,KB
region. It then walks a binary tree of subregions upward from the faulting
page and, at each level, compares the count of set pages with the
subregion's size. The largest subregion for which
\begin{equation}
  100 \cdot \textit{count} \;>\; \textit{threshold} \cdot \textit{size}
  \label{eq:density}
\end{equation}
becomes the prefetch region. Its non-resident pages are added to the batch's
destination set and marked as prefetches.

Three details of this rule drive Sections~VI and~VII.

\paragraph{The rule counts residency, not mappings}
The bitmap is the destination's resident mask combined with the faulted
pages; mapping state does not enter it. A page made resident on the GPU by
any means, including a mechanism outside the fault path, counts toward the
density test exactly as a page brought in by an earlier fault does.

\paragraph{Prefetched pages are mapped in the same step}
The prefetch region is merged into the set of pages being serviced, and the
mapping loop at the end of block service runs over that whole set, including
pages that needed no copy because they were already resident. A region that
passes the test is therefore mapped for the GPU during the fault that
triggered it, which prevents the later faults on its other pages.

\paragraph{The threshold is a load-time parameter}
The threshold is the module parameter \texttt{uvm\_perf\_prefetch\_threshold},
a percentage with default 51, set when the module is loaded. The test in
(\ref{eq:density}) is strict. At 100 it can never pass, which disables the
density rule; at 0 it passes for any subregion containing a single counted
page, so the whole region the prefetcher considers is taken on the first
fault in it. Values above 100 are rejected silently: the driver falls back to
51 while the parameter file still reports the value that was written, a
hazard our experiments guard against by keeping every setting within range
and reading back the effective behaviour rather than the parameter.

A second prefetch path populates an entire block on its first touch, but only
when the block is resident nowhere and the faulting GPU is the block's
preferred location. Our benchmarks initialise their data on the host and set
no preferred location, so this path does not fire in our measurements; it
can fire for applications that issue \texttt{cudaMemAdvise}.

\subsection{Prefetch Aggressiveness Under Memory Pressure}
\label{sec:bg-oversub}

When GPU memory is full, servicing a fault requires evicting another block's
data to the host, and the driver selects victims at block granularity
\cite{go2023earlyadaptor,ganguly2019interplay}. The density threshold
therefore trades one cost against another. A low threshold migrates large
regions early, which reduces the number of faults when most of a region will
be used. When only a few pages of each region are used, the same setting
moves mostly data that is never touched; if the inflated footprint exceeds
GPU memory, it also evicts data that will be needed again. Go
\emph{et al.}\ measured both sides of this trade-off on an RTX~3090 and built
an adaptive threshold controller on it \cite{go2023earlyadaptor}. Section~VII
measures it on newer hardware and driver releases, and Section~VI shows that
off-path speculation, as we built it, acts on performance only through this
same rule.

\subsection{Where Off-Path Work Can Act}
\label{sec:bg-where}

Combining the three subsections gives the frame for the rest of the paper.
Work moved off the fault path, by a background thread that predicts upcoming
faults, can touch only residency. It can therefore reduce the copy inside D5,
and it can raise the count in (\ref{eq:density}) so that the in-path
prefetcher maps a region sooner than it otherwise would. It cannot install a
mapping itself, and it does not shorten the lock-held classification and
bookkeeping that make up most of D5. Whatever it gains has to arrive through
one of those two channels, and both are measurable: the first as a change in
D5 per fault, the second as a change in the number of demand faults at a
fixed threshold.

% [FIG] Fig. 1: one fault batch D1-D6 on a timeline; inside D5, the steps
% classify -> residency -> prefetch hint (Eq. 1) -> copy -> map; arrows marking
% where (a) a resident-only stager acts (copy, Eq. 1 count) and (b) an explicit
% migrate-and-map call acts (residency + mapping, outside the fault path).
